\chapter{Минимальная опора кривизны и no-go единственного общего следа}
\label{ch:v7-minimal-curvature-support-trace-gate}

\section{От контейнера к структурной опоре}

Предыдущий гейт доказал, что средняя ступень связывающей цепи
тождественно аннулируется relative curvature. Поэтому допустимое сжатие
имеет вид
\begin{equation}
 P_{02}:H_0\oplus H_1\oplus H_2\longrightarrow H_0\oplus H_2,
 \qquad \dim(H_0\oplus H_2)=11+10=21.
\end{equation}
На восемнадцати физических конфигурациях проверено
\begin{equation}
 \Tr_{42}\mathbb M_L^2
 =\Tr_{21}(P_{02}\mathbb M_LP_{02})^2
 \label{eq:v7-support-trace-compression}
\end{equation}
с максимальной невязкой $4.6\times10^{-13}$.

Нельзя сжимать дальше к поточечному образу самой кривизны. На текущем
физическом срезе её общий ранг достигает $14$, но в начале равен нулю.
Норма скачка опорного проектора от нулевой точки к общей ненулевой точке
равна
\begin{equation}
 \|P_{\operatorname{supp}\mathbb M_L}-0\|_{\rm HS}=\sqrt{14}.
 \label{eq:v7-support-projector-jump}
\end{equation}
Следовательно, нормировка по поточечному рангу не задаёт фиксированный
гильбертов носитель и не является допустимой заменой одного следа.

\section{Порождённая алгебра связывающего сектора}

Для сжатых матриц решена система
\begin{equation}
 [X,P_{02}\mathbb M_L(A)P_{02}]=0
\end{equation}
на восемнадцати общих физических конфигурациях. Коммутант имеет комплексную
размерность один; наименьшее ненулевое сингулярное число системы равно
$5.2050935227\ldots$. По конечномерной теореме Бернсайда порождённая
${}^*$-алгебра неприводима:
\begin{equation}
 \mathcal A_L=M_{21}(\mathbb C).
 \label{eq:v7-support-linking-factor}
\end{equation}
На этом одном факторе нормированный след действительно единственен.

\section{Почему рёберный суррогат не является фактором}

Рёберная кривизна общего 96-мерного контейнера была представлена как
\begin{equation}
 \mathbb M_E=\operatorname{diag}(m_E,-m_E),
 \qquad m_E\in\mathbb R^{27}.
\end{equation}
Двадцать семь независимых знаковых диагональных генераторов и их квадраты
порождают все $54$ диагональные матричные единицы. Поэтому
\begin{equation}
 \mathcal A_E\simeq\mathbb C^{54},
 \qquad
 \dim\mathcal A_E=\dim\mathcal A_E'=54.
 \label{eq:v7-support-edge-abelian-algebra}
\end{equation}
Это точный следовой суррогат полного гессиана, но не простой физический
фактор.

В результате фактически порождённая алгебра равна
\begin{equation}
 \mathcal A_{\rm gen}=\mathbb C^{54}\oplus M_{21}(\mathbb C).
 \label{eq:v7-support-generated-algebra}
\end{equation}
Её центр имеет размерность $55$. После нормировки верный положительный след
содержит $54$ независимых параметра. Уже выведенные Casimir- и
Hodge-веса фиксируют внутреннее рёберное действие, но между полным рёберным
и связывающим функционалами всё равно остаётся как минимум один
относительный вес.

Если вместо диагонального суррогата вернуться к исходному физическому
Hodge-носителю, получается уже проверенная развилка
\begin{equation}
 M_{22}(\mathbb C)\oplus M_{21}(\mathbb C),
 \label{eq:v7-support-physical-factor-fork}
\end{equation}
центр которой двумерен и оставляет один относительный следовой параметр.

\section{Почему полный матричный след нельзя объявить}

Стандартный след одной копии общей матрицы размера $75$ ограничивается на
два угла с весами
\begin{equation}
 \frac{54}{75}=\frac{18}{25},
 \qquad
 \frac{21}{75}=\frac7{25},
\end{equation}
и воспроизводит равный ненормированный коэффициент двух блоков. Но это
свойство выбранного представления контейнера, а не абстрактной прямой суммы.
Чтобы алгебра действительно стала простой $M_{75}(\mathbb C)$, необходимо
добавить
\begin{equation}
 \operatorname{Hom}(\mathbb C^{21},\mathbb C^{54}),
 \qquad \dim_{\mathbb C}=54\cdot21=1134
 \label{eq:v7-support-m75-connector-cost}
\end{equation}
вместе с сопряжённым углом. Текущий проект не содержит канонического
оператора этого типа.

Аналогично физическое объединение $M_{22}\oplus M_{21}$ в простой
$M_{43}$ требует $22\cdot21=462$ комплексных компонент. Этот маршрут уже
был закрыт гейтом кратностей общего неприводимого следа.

Именно так различаются положительные прецеденты Томов V--VI и текущий
контейнер. В прежних случаях полная простая алгебра, переходные операторы и
нецентральные угловые проекторы были построены явно. Классификация конечных
спектральных троек также требует представления и матриц кратностей, а не
одного общего размера \cite{v7-support-krajewski,v7-support-cacic,
v7-support-masson}.

\section{Контроль качественного вакуума}

Для относительных весов
\begin{equation}
 \eta=\frac14,\ 1,\ \frac{18}{7},\ 4
\end{equation}
полный 27-мерный гессиан сохраняет
\begin{equation}
 (n_-,n_0,n_+)_{0}=(7,0,20),
 \qquad
 (n_-,n_0,n_+)_{*}=(0,0,27).
\end{equation}
Но минимальная вакуумная мода меняется от $3.7684\ldots$ до
$5.4286\ldots$, а максимальная --- от $24$ до $69.3205\ldots$.
Качественное замыкание сохраняется, количественный спектр не фиксирован.

\begin{theorem}[No-go следа минимальной опоры]
Сжатие к фиксированной концевой опоре размерности $21$ канонично и
сохраняет связывающее действие. Однако порождённая общая алгебра остаётся
прямой суммой с нетривиальным центром. Ни верность, ни нормировка следа не
устраняют относительный центральный вес. Объявление простого полного
матричного фактора требует нового внедиагонального коннектора.
\end{theorem}

\begin{remark}
Тем самым исчерпаны размерностный, кратностный, Real-, клиффордов,
бикомплексный и опорный маршруты количественной нормировки текущего
родителя. Ни один из них не затрагивает нормировочно независимый селектор и
строго устойчивый локальный вакуум.
\end{remark}

\begin{nogocriterion}[Запрет простоты контейнера]
Нельзя считать блочную запись в $M_{75}$ простой координатной алгеброй до
предъявления ненулевого физического блока между $\mathbb C^{54}$ и
$\mathbb C^{21}$. Также нельзя нормировать след на поточечный ранг
кривизны, поскольку её опорный проектор скачкообразно меняется в начале.
\end{nogocriterion}

\begin{protocol}
\begin{enumerate}
 \item структурная связывающая опора: \textbf{$21$};
 \item максимальный поточечный ранг: \textbf{$14$};
 \item поточечная ранговая нормировка: \textbf{закрыта};
 \item коммутант связывающего семейства: \textbf{размерность $1$};
 \item связывающая алгебра: \textbf{$M_{21}(\mathbb C)$};
 \item рёберная суррогатная алгебра: \textbf{$\mathbb C^{54}$};
 \item центр общей порождённой алгебры: \textbf{размерность $55$};
 \item свободных нормированных следовых параметров: \textbf{$54$};
 \item физическая факторная модель: \textbf{$M_{22}\oplus M_{21}$};
 \item её относительных весов: \textbf{$1$};
 \item цена простого $M_{75}$: \textbf{$1134$ комплексных коннектора};
 \item цена простого $M_{43}$: \textbf{$462$ комплексных коннектора};
 \item канонический коннектор: \textbf{не найден};
 \item качественный родитель: \textbf{замкнут};
 \item массовая метрика: \textbf{не выведена};
 \item следующий гейт: \textbf{заморозка качественного родителя и границы масс}.
\end{enumerate}
\end{protocol}

Машинный сертификат сохранён в
\path{s2t_v7_minimal_curvature_support_trace_gate_results.json}.

\begin{thebibliography}{9}
\bibitem{v7-support-krajewski}
T.~Krajewski,
\emph{Classification of Finite Spectral Triples},
J. Geom. Phys. 28 (1998), 1--30; arXiv:hep-th/9701081.
\bibitem{v7-support-cacic}
B.~\v{C}a\'ci\'c,
\emph{Moduli Spaces of Dirac Operators for Finite Spectral Triples},
J. Noncommut. Geom. 3 (2009), 453--495; arXiv:0902.2068.
\bibitem{v7-support-masson}
T.~Masson, G.~Nieuviarts,
\emph{Lifting Bratteli Diagrams between Krajewski Diagrams: Spectral
Triples, Spectral Actions, and AF Algebras}, arXiv:2207.04466.
\end{thebibliography}